% chapters/ch06_spi_i2c.tex — Chapter 6: SPI & I2C
% Scope (PLAN.md): timing diagrams, 4 SPI modes (CPOL/CPHA), chip select,
% I2C start/stop/ACK, addressing, clock stretching, arbitration, pull-up
% sizing, multi-slave topologies, SPI vs I2C vs UART selection.
% Budget 30 pp, mix 18 Q&A / 9 code / 3 concept.
% All snippets verified gcc -Wall -Wextra -std=c11 (zero warnings).

\chapter{SPI \& I2C}
\label{ch:spi-i2c}

\begin{partnote}
\textbf{Scope} --- SPI clock polarity/phase (the four modes) $\rightarrow$
chip-select and multi-slave topologies $\rightarrow$ I2C start/stop/ACK
$\rightarrow$ 7- and 10-bit addressing $\rightarrow$ clock stretching
$\rightarrow$ arbitration $\rightarrow$ pull-up sizing $\rightarrow$ choosing
SPI vs I2C vs UART. These are the two short-range \emph{synchronous} buses
(a shared clock), the counterpart to Chapter 5's asynchronous UART.
\end{partnote}

SPI and I2C both carry a clock alongside the data, so unlike UART there is no
baud-matching --- the master clocks the slave. They trade differently: SPI is
fast, simple, and pin-hungry; I2C is two-wire, addressable, and slower. Knowing
\emph{which} to pick, and the timing rules that make each work, is the core of
this chapter.

%=====================================================================
\section{Primer}
%=====================================================================

\subsection{SPI: full-duplex, four modes, one select per slave}

SPI is a synchronous, full-duplex, master-driven bus with four lines:
\textbf{SCLK} (clock), \textbf{MOSI} (master-out), \textbf{MISO} (master-in),
and one \textbf{CS/SS} (chip-select, active-low) \emph{per slave}. On each clock
the master shifts a bit out on MOSI and simultaneously shifts one in on MISO ---
data goes both ways at once. Two settings define the timing: \textbf{CPOL}
(clock idle level, 0 = idle low) and \textbf{CPHA} (which clock edge samples
data: 0 = first edge, 1 = second edge). The four combinations are SPI
\textbf{modes 0--3}; master and slave must use the same mode or every bit is
mis-sampled. There's no addressing and no ACK --- selection is purely the CS
line, so $N$ slaves need $N$ chip-selects (or a daisy-chain).

\subsection{I2C: two wires, addressing, open-drain}

I2C is a synchronous, half-duplex, two-wire bus: \textbf{SDA} (data) and
\textbf{SCL} (clock), both \textbf{open-drain} with pull-up resistors (Chapter
3), so any device can pull low and release. A transaction is framed by a
\textbf{START} (SDA falls while SCL is high) and a \textbf{STOP} (SDA rises
while SCL is high). The master sends a 7-bit (or 10-bit) \textbf{address} plus
a read/write bit; the addressed slave \textbf{ACK}s by pulling SDA low on the
9th clock. Every byte is acknowledged, so I2C has built-in delivery
confirmation UART and SPI lack. Because it's open-drain and addressed, many
devices share just two wires --- the trade is speed (standard 100\,kHz, fast
400\,kHz, vs SPI's tens of MHz) and protocol complexity.

\subsection{Clock stretching, arbitration, and pull-ups}

Two I2C mechanisms matter. \textbf{Clock stretching}: a slow slave holds SCL low
to pause the master until it's ready --- the master must tolerate this (and time
out if a slave hangs SCL forever). \textbf{Arbitration}: if two masters start
together, each watches SDA while driving; the one that wanted to send a 1 but
sees 0 (another master pulled it low) loses and backs off, non-destructively ---
the wired-AND bus makes this work. \textbf{Pull-up sizing}: the resistors and
the bus capacitance form an RC that sets the rising-edge time; too large a
resistor and the edge is too slow for 400\,kHz, too small and the open-drain
sink current is excessive --- a few k$\Omega$ is typical, computed from the
spec's rise-time limit and total bus capacitance.

%=====================================================================
\section{Q\&A Bank}
%=====================================================================

\tierbanner{Tier 1 --- Screening (phone / HR-technical)}%
{Two to four sentences.}

\question{Q1.1}{Name the SPI signals and what each does.}
SCLK (clock from the master), MOSI (master-out slave-in), MISO (master-in
slave-out), and one active-low CS/SS per slave. Each clock edge shifts a bit out
on MOSI and in on MISO simultaneously --- SPI is full-duplex.

\question{Q1.2}{Name the I2C signals and why they're open-drain.}
SDA (data) and SCL (clock), both open-drain with pull-ups so any device can pull
the line low or release it (high-Z). This wired-AND arrangement lets multiple
devices share two wires without driver contention.

\question{Q1.3}{What are CPOL and CPHA?}
CPOL sets the clock's idle level (0 = idle low, 1 = idle high); CPHA selects
which clock edge samples data (0 = first edge, 1 = second). Their four
combinations are SPI modes 0--3.

\question{Q1.4}{How does a master select a slave on SPI vs I2C?}
SPI: assert that slave's dedicated chip-select line (one per slave). I2C: send
the slave's address on the shared bus and the matching slave responds --- no
extra wires per device.

\question{Q1.5}{What is an I2C ACK?}
After each byte the receiver pulls SDA low on the 9th clock to acknowledge
receipt; leaving it high is a NACK (not acknowledged) --- used to signal ``no
such device,'' an error, or ``last byte, stop.''

\question{Q1.6}{Which is faster, SPI or I2C, and which uses fewer pins?}
SPI is faster (tens of MHz, full-duplex) but needs four lines plus one CS per
slave. I2C is slower (100\,k--400\,kHz typically) but needs only two wires for
many devices.

\question{Q1.7}{What frames an I2C transaction?}
A START condition (SDA falling while SCL high) begins it and a STOP (SDA rising
while SCL high) ends it. Between them are address and data bytes, each
ACK/NACKed.

\question{Q1.8}{Is SPI addressed?}
No --- SPI has no addresses or ACKs; the master picks a slave purely by
asserting its chip-select. Selection is hardware, not protocol.

\tierbanner{Tier 2 --- Core technical (panel round)}%
{A short paragraph.}

\question{Q2.1}{Walk through the four SPI modes and the consequence of a
mismatch.}
Mode = (CPOL, CPHA): mode 0 (0,0) idle-low, sample on the rising (first) edge;
mode 1 (0,1) idle-low, sample on falling (second); mode 2 (1,0) idle-high,
sample on falling (first); mode 3 (1,1) idle-high, sample on rising (second).
Master and slave must agree, because CPHA decides whether data is sampled on the
edge that the other side is \emph{changing} it --- a mismatch samples mid-
transition and every bit is wrong or shifted by half a clock. The datasheet of
the slave fixes the mode; you configure the master to match.

\question{Q2.2}{How do you connect multiple SPI slaves?}
Two ways. \textbf{Independent CS}: shared SCLK/MOSI/MISO, one chip-select line
per slave --- simple, but costs a GPIO per device and only one slave is selected
at a time. \textbf{Daisy-chain}: slaves' MISO$\rightarrow$MOSI are chained so
data shifts through all of them with one CS --- saves pins but needs device
support and shifting the full chain length. Independent CS is by far the common
choice; watch that unselected slaves tri-state MISO so they don't fight the bus.

\question{Q2.3}{Explain the full I2C write transaction byte by byte.}
START $\rightarrow$ address byte (7-bit address + W=0) $\rightarrow$ slave ACK
$\rightarrow$ register/command byte $\rightarrow$ ACK $\rightarrow$ data byte(s)
each followed by ACK $\rightarrow$ STOP. A read is START $\rightarrow$ address+W
$\rightarrow$ register $\rightarrow$ repeated START $\rightarrow$ address+R
$\rightarrow$ slave sends bytes, master ACKs each except NACKs the last
$\rightarrow$ STOP. The \emph{repeated start} (no STOP between) keeps the bus so
no other master interrupts the read.

\question{Q2.4}{What is clock stretching and how must the master handle it?}
A slave that needs more time holds SCL low after an ACK, pausing the clock until
it releases. The master must \emph{not} assume it owns the clock timing --- it
samples whether SCL actually rose before proceeding, and applies a timeout so a
slave that hangs SCL low forever doesn't lock the bus indefinitely (a common
recovery is clocking SCL manually to free a stuck slave). Hardware I2C
peripherals handle stretching automatically; bit-banged ones must poll SCL.

\question{Q2.5}{How does I2C multi-master arbitration work?}
All masters can start; each drives SDA and simultaneously reads it back. Because
the bus is wired-AND, a master that tries to send a 1 but reads 0 knows another
master pulled it low --- it has lost arbitration and immediately stops driving,
non-destructively (the winning master's transfer is unaffected and continues).
Arbitration is decided bit by bit on the address/data, so the lower-address (or
the one sending more zeros first) wins. No data is corrupted.

\question{Q2.6}{How do you size I2C pull-up resistors?}
The pull-up and the total bus capacitance form an RC; the resistor must be small
enough that the rising edge reaches the high threshold within the spec's
rise-time limit (1000\,ns for standard mode, 300\,ns for fast mode), and large
enough that the low-side sink current (V/R) stays within the open-drain output's
rating (3\,mA). So $R_{\min}=V_{cc}/I_{\text{sink}}$ and $R_{\max}$ from the
rise-time/capacitance limit --- typically landing a few k$\Omega$ (e.g.\
4.7\,k$\Omega$ at 100\,kHz, 2.2\,k$\Omega$ at 400\,kHz). More devices = more
capacitance = smaller resistor needed.

\question{Q2.7}{Compute the SPI clock from the prescaler.}
On most MCUs SPI SCLK = $f_{\text{PCLK}} / 2^{(\text{presc}+1)}$, i.e.\
prescaler values select divide-by-2, 4, 8, \ldots\ So an 84\,MHz PCLK with
prescaler index 2 gives $84/8 = 10.5$\,MHz. You pick the highest divider whose
result stays within the slave's maximum SCLK \emph{and} within the board's
signal-integrity limit (long traces ring at high SCLK). Unlike UART there's no
``error'' --- the slave is clocked by the master, so any rate up to its max
works.

\question{Q2.8}{When do you choose SPI vs I2C vs UART?}
\textbf{UART} for asynchronous, point-to-point, no-shared-clock links
(consoles, modules, RS422). \textbf{SPI} when you need speed and full-duplex and
can spare the pins --- ADCs, displays, flash, radios. \textbf{I2C} when you have
many slow peripherals and want minimal wiring --- sensors, EEPROMs, config
chips. The quick decision: async peer $\rightarrow$ UART; high throughput
$\rightarrow$ SPI; pin-constrained many-device bus $\rightarrow$ I2C.

\question{Q2.9}{Why does I2C have an inherent speed ceiling that SPI doesn't?}
Open-drain drive can only \emph{pull} low; the rising edge is restored passively
by the pull-up through the bus capacitance, so the RC rise time bounds the clock
--- you can't go faster than the line can rise. SPI is push-pull (actively
drives both edges) and point-to-point with short, defined loading, so it clocks
much faster. That same open-drain property is what gives I2C its multi-device,
wired-AND, arbitration-friendly nature --- speed is the trade for
shareability.

\tierbanner{Tier 3 --- Trap \& depth (principal-engineer level)}%
{A paragraph plus the trap.}

\question{Q3.1}{Your SPI reads are consistently shifted by one bit or return
garbage on a new slave. What do you check, in order?}
First the \textbf{mode}: the slave's datasheet CPOL/CPHA must match the master
--- a CPHA mismatch samples on the changing edge and shifts/garbles every bit.
Then \textbf{bit order} (MSB-first vs LSB-first), \textbf{word size} (8 vs 16
bit), and \textbf{CS timing} (some slaves require CS to frame each word, or a
setup time before the first clock). Then signal integrity: too-high SCLK on long
traces rings and the slave mis-samples. \trap the reflexive ``lower the clock''
sometimes helps signal integrity and masks a real mode mismatch --- the
disciplined order is mode $\rightarrow$ bit order $\rightarrow$ word size
$\rightarrow$ CS framing $\rightarrow$ scope the SCLK/MOSI/MISO at the slave.
This is the same scope-at-the-receiver method as the RS422 chapter.

\question{Q3.2}{An I2C bus is ``stuck'' --- SDA or SCL held low and no
transactions complete. Diagnose and recover.}
A slave can hold SDA low if it was mid-byte when the master was reset (it's
waiting to finish clocking out a byte), or hold SCL low via clock stretching
that never releases (a hung slave). Recovery: the master manually toggles SCL up
to 9 clocks to let the slave finish its byte and release SDA, then issues a STOP
to resync the bus; if SCL itself is stuck, the slave or a bus buffer has failed
and only a power cycle / reset line clears it. \trap candidates try to ``just
re-init the I2C peripheral,'' which doesn't move SCL and so doesn't free a slave
mid-byte. The correct recovery is the manual 9-clock bus-clear sequence ---
knowing it signals real I2C field experience.

\question{Q3.3}{Why can a single missing or wrong-valued I2C pull-up cause
intermittent, capacitance-dependent failures?}
The pull-up sets the rising-edge time through the bus RC; if it's too weak (too
large) or absent (relying on a weak internal pull-up), the edge is too slow and
the receiver samples before the line crosses the high threshold --- worse as you
add devices (more capacitance) or raise the clock to 400\,kHz, so it ``works on
the bench, fails on the full board.'' Too strong (too small) overloads the
open-drain low-side and may not reach a valid low. \trap the trap is treating
pull-ups as ``any value works'' --- they're a computed RC trade against bus
capacitance and clock, and a marginal value is a textbook source of
intermittent, layout-dependent I2C bugs that a scope (slow rise time) reveals
instantly.

\question{Q3.4}{SPI has no ACK and no addressing --- how do you build a reliable
multi-slave SPI system, and what failure modes must you guard?}
Because SPI gives no delivery confirmation, reliability is layered on top:
per-transaction framing with a command/length/CRC so the master can verify the
slave's response, a known ``idle/dummy'' pattern to detect an absent or wedged
slave (all-0xFF or all-0x00 reads usually mean MISO floating / slave not
driving), and careful CS sequencing so exactly one slave drives MISO at a time
(an unselected slave that doesn't tri-state corrupts the bus). \trap the trap is
trusting raw SPI bytes as if they were acknowledged --- they aren't, so a
disconnected or mis-clocked slave returns plausible-looking garbage. The robust
design validates a CRC/known-pattern in every response and treats all-ones/all-
zeros as ``no slave,'' exactly the frame-integrity discipline from Chapter 5.

%=====================================================================
\section{Coding Gym}
%=====================================================================

\begin{partnote}
Nine host-compilable problems (verified under
\code{gcc -Wall -Wextra -std=c11}) covering the bit/protocol arithmetic of SPI
and I2C bring-up: mode decode, software transfer, clock dividers, address
encoding, ACK logic, clock-stretch handling, and bus selection.
\end{partnote}

%---------------------------------------------------------------
\begin{gymproblem}{Decode an SPI mode to CPOL/CPHA}

\textbf{Problem.} Convert an SPI mode number (0--3) into its CPOL and CPHA
bits.

\textbf{Hints.}
\begin{itemize}
\item Mode = (CPOL$<<$1) $|$ CPHA.
\end{itemize}

\attempt

\textbf{Solution.}
\begin{cgym}
#include <stdio.h>

static void spi_mode(int mode, int *cpol, int *cpha) {
    *cpol = (mode >> 1) & 1;
    *cpha = mode & 1;
}

int main(void) {
    int cpol, cpha;
    spi_mode(3, &cpol, &cpha);
    printf("mode3 cpol=%d cpha=%d\n", cpol, cpha);   /* 1 1 */
    return 0;
}
\end{cgym}

Mode 3 is CPOL=1 (idle high), CPHA=1 (sample second edge). The first thing you
do bringing up an SPI slave is read its datasheet timing diagram, derive the
mode, and set the master to match --- the \#1 cause of garbled SPI is a mode
mismatch (Q3.1).

\followups
\begin{enumerate}
\item \emph{``Build the reverse: mode from CPOL/CPHA.''} $\rightarrow$
  \code{(cpol<<1)|cpha}.
\item \emph{``Which edge does mode 0 sample?''} $\rightarrow$ Rising (first)
  edge, clock idle low.
\item \emph{``Two slaves need different modes on one bus --- handling?''}
  $\rightarrow$ Reconfigure the master's mode per CS, or use separate SPI
  peripherals.
\end{enumerate}
\end{gymproblem}

%---------------------------------------------------------------
\begin{gymproblem}{Software SPI byte transfer (MSB-first)}

\textbf{Problem.} Model a full-duplex 8-bit SPI exchange: shift a byte out
while shifting one in, MSB-first.

\textbf{Hints.}
\begin{itemize}
\item Loop bits 7..0; build the input by shifting in the slave's bit each clock.
\end{itemize}

\attempt

\textbf{Solution.}
\begin{cgym}
#include <stdint.h>
#include <stdio.h>

/* models the simultaneous shift: returns the byte clocked in from the slave */
static uint8_t spi_xfer_mode0(uint8_t mosi, uint8_t slave_shift) {
    uint8_t miso = 0;
    for (int i = 7; i >= 0; i--) {        /* MSB first */
        int out = (mosi >> i) & 1; (void)out;   /* would drive MOSI here */
        int in  = (slave_shift >> i) & 1;       /* sample MISO on the edge */
        miso = (uint8_t)((miso << 1) | in);
    }
    return miso;
}

int main(void) {
    printf("in=0x%02X\n", spi_xfer_mode0(0xFF, 0xA5));   /* 0xA5 */
    return 0;
}
\end{cgym}

The defining SPI property is here: \emph{every clock both transmits and
receives}. A real bit-bang drives MOSI then toggles SCLK then samples MISO at
the CPHA-defined edge; a hardware SPI does it in one register write. MSB-first
is the common default --- LSB-first slaves exist and a wrong bit order reverses
every byte.

\followups
\begin{enumerate}
\item \emph{``Add CPHA: when do you sample vs change?''} $\rightarrow$ CPHA=0
  sample on the first edge (change before SCLK starts); CPHA=1 sample on the
  second.
\item \emph{``A read-only register --- what do you send?''} $\rightarrow$ Dummy
  bytes (0x00/0xFF) to generate clocks; the data comes back on MISO.
\item \emph{``16-bit transfers?''} $\rightarrow$ Same loop to 16; or two 8-bit
  transfers with CS held.
\end{enumerate}
\end{gymproblem}

%---------------------------------------------------------------
\begin{gymproblem}{SPI clock from the prescaler}

\textbf{Problem.} Compute SCLK from the peripheral clock and a power-of-two
prescaler index.

\textbf{Hints.}
\begin{itemize}
\item SCLK = PCLK / $2^{(\text{presc}+1)}$.
\end{itemize}

\attempt

\textbf{Solution.}
\begin{cgym}
#include <stdint.h>
#include <stdio.h>

static uint32_t spi_sclk(uint32_t pclk, int presc) {
    return pclk >> (presc + 1);     /* /2, /4, /8, ... */
}

int main(void) {
    printf("sclk=%u\n", spi_sclk(84000000u, 2));   /* 84MHz/8 = 10.5 MHz */
    return 0;
}
\end{cgym}

Prescaler index 2 means divide-by-8, giving 10.5\,MHz. You pick the fastest rate
within the slave's max SCLK and the board's signal integrity. Unlike UART there
is no baud error to worry about --- the master clocks the slave directly --- but
too-high SCLK on long traces causes ringing and mis-sampling, the SPI analogue
of an unterminated line.

\followups
\begin{enumerate}
\item \emph{``Slave max is 8\,MHz --- which prescaler from 84\,MHz?''}
  $\rightarrow$ presc=3 ($/16$ = 5.25\,MHz); presc=2 ($/8$ = 10.5) exceeds it.
\item \emph{``Why no baud error like UART?''} $\rightarrow$ The clock is shared,
  so there's no sampling drift between ends.
\item \emph{``Faster than the slave's max --- symptom?''} $\rightarrow$
  Corrupt/garbled reads; back off the prescaler.
\end{enumerate}
\end{gymproblem}

%---------------------------------------------------------------
\begin{gymproblem}{I2C 7-bit address byte}

\textbf{Problem.} Form the first I2C byte from a 7-bit address and the
read/write direction.

\textbf{Hints.}
\begin{itemize}
\item Address in bits 7..1, R/W in bit 0 (1 = read).
\end{itemize}

\attempt

\textbf{Solution.}
\begin{cgym}
#include <stdint.h>
#include <stdio.h>

static uint8_t i2c_addr_byte(uint8_t addr7, int read) {
    return (uint8_t)((addr7 << 1) | (read ? 1 : 0));
}

int main(void) {
    printf("read=0x%02X write=0x%02X\n",
           i2c_addr_byte(0x48, 1), i2c_addr_byte(0x48, 0));   /* 0x91 0x90 */
    return 0;
}
\end{cgym}

Address 0x48 (a common temperature-sensor address) becomes 0x90 for a write and
0x91 for a read. A frequent confusion is the ``8-bit vs 7-bit address''
mismatch: datasheets sometimes quote the already-shifted 8-bit value (0x90),
sometimes the 7-bit (0x48) --- mixing them addresses the wrong device. Always
know which form your API expects.

\followups
\begin{enumerate}
\item \emph{``Datasheet says address 0x90 --- what's the 7-bit value?''}
  $\rightarrow$ \code{0x90 >> 1 = 0x48}.
\item \emph{``Reserved addresses?''} $\rightarrow$ 0x00--0x07 and 0x78--0x7F are
  reserved (general call, 10-bit, etc.).
\item \emph{``Scan the bus for devices.''} $\rightarrow$ Address each 0x08--0x77
  and watch for an ACK.
\end{enumerate}
\end{gymproblem}

%---------------------------------------------------------------
\begin{gymproblem}{I2C 10-bit addressing}

\textbf{Problem.} Build the two address bytes for 10-bit I2C addressing.

\textbf{Hints.}
\begin{itemize}
\item First byte: \code{11110} prefix + top 2 address bits + R/W.
\item Second byte: the low 8 address bits.
\end{itemize}

\attempt

\textbf{Solution.}
\begin{cgym}
#include <stdint.h>
#include <stdio.h>

static void i2c_addr10(uint16_t addr10, int read, uint8_t *b1, uint8_t *b2) {
    *b1 = (uint8_t)(0xF0 | (((addr10 >> 8) & 0x3) << 1) | (read ? 1 : 0));
    *b2 = (uint8_t)(addr10 & 0xFF);
}

int main(void) {
    uint8_t b1, b2;
    i2c_addr10(0x2A5, 0, &b1, &b2);
    printf("b1=0x%02X b2=0x%02X\n", b1, b2);   /* 0xF4 0xA5 */
    return 0;
}
\end{cgym}

The \code{11110xx} prefix tells slaves a 10-bit address follows; the top two
address bits ride in the first byte, the low eight in the second. 10-bit
addressing is rare (most devices are 7-bit) but appears when a bus needs more
than 112 usable addresses; recognising the \code{0xF0} prefix on a logic
analyser is the tell.

\followups
\begin{enumerate}
\item \emph{``Why is 7-bit far more common?''} $\rightarrow$ 112 usable
  addresses suffice for most boards, with less protocol overhead.
\item \emph{``A 10-bit read uses how many phases?''} $\rightarrow$ Write the two
  address bytes, repeated START, then the first byte again with R=1.
\item \emph{``Mixed 7- and 10-bit on one bus?''} $\rightarrow$ Allowed --- the
  \code{11110} prefix is reserved so 7-bit devices ignore it.
\end{enumerate}
\end{gymproblem}

%---------------------------------------------------------------
\begin{gymproblem}{I2C SCL frequency from timing registers}

\textbf{Problem.} Compute the SCL clock from the peripheral clock, prescaler,
and the low/high period counts.

\textbf{Hints.}
\begin{itemize}
\item Period = (SCLL+1) + (SCLH+1) prescaled ticks; SCL = PCLK / (presc $\times$
  period).
\end{itemize}

\attempt

\textbf{Solution.}
\begin{cgym}
#include <stdint.h>
#include <stdio.h>

static uint32_t i2c_scl_hz(uint32_t pclk, uint32_t presc, uint32_t scll, uint32_t sclh) {
    uint32_t tpresc = presc + 1u;
    return pclk / (tpresc * ((scll + 1u) + (sclh + 1u)));
}

int main(void) {
    printf("scl=%u Hz\n", i2c_scl_hz(8000000u, 1u, 0x9u, 0x9u));   /* 200000 */
    return 0;
}
\end{cgym}

The SCL period is the programmed low time plus high time (each in prescaled
ticks); here 8\,MHz with prescaler 2 and SCLL=SCLH=9 gives
$8\text{M}/(2\times20) = 200$\,kHz. Standard mode is 100\,kHz, fast mode
400\,kHz --- and the high/low split isn't 50/50 in the spec (the low period is
longer), which is why real timing registers separate SCLL and SCLH. Set them
from the datasheet's $t_{\text{LOW}}/t_{\text{HIGH}}$ minimums.

\followups
\begin{enumerate}
\item \emph{``Why is the low period longer than the high?''} $\rightarrow$ The
  spec's $t_{\text{LOW}} > t_{\text{HIGH}}$ to allow slow rising edges and
  setup.
\item \emph{``400\,kHz from 8\,MHz --- pick values.''} $\rightarrow$ Reduce the
  period count so PCLK/(presc$\times$period)=400k.
\item \emph{``Rise time eats into the high period --- effect?''} $\rightarrow$
  Effective SCL high shrinks; weak pull-ups can violate timing.
\end{enumerate}
\end{gymproblem}

%---------------------------------------------------------------
\begin{gymproblem}{Master ACK/NACK logic on a multi-byte read}

\textbf{Problem.} On an I2C read of \code{total} bytes, decide whether the
master ACKs after byte \code{i}.

\textbf{Hints.}
\begin{itemize}
\item ACK every byte except the last; NACK the last to tell the slave to stop.
\end{itemize}

\attempt

\textbf{Solution.}
\begin{cgym}
#include <stdio.h>

static int master_should_ack(int byte_index, int total) {
    return byte_index < (total - 1);     /* NACK only the final byte */
}

int main(void) {
    printf("ack[0/3]=%d ack[2/3]=%d\n",
           master_should_ack(0, 3), master_should_ack(2, 3));   /* 1 0 */
    return 0;
}
\end{cgym}

During a read the master ACKs each received byte to say ``send the next,'' then
\textbf{NACKs} the last one to tell the slave ``I'm done, release SDA'' before
issuing STOP. Forgetting the final NACK is a classic bug: the slave keeps
driving and you either read an extra byte or the STOP collides with the slave
still holding SDA --- a wedged bus (Q3.2).

\followups
\begin{enumerate}
\item \emph{``What if you ACK the last byte too?''} $\rightarrow$ The slave
  prepares another byte and may hold the bus; STOP can be mis-issued.
\item \emph{``Reading an unknown-length stream?''} $\rightarrow$ Some devices
  use a length byte first, or you NACK when your buffer is full.
\item \emph{``Who ACKs on a write?''} $\rightarrow$ The \emph{slave} ACKs each
  byte the master sends.
\end{enumerate}
\end{gymproblem}

%---------------------------------------------------------------
\begin{gymproblem}{Wait out clock stretching (with timeout)}

\textbf{Problem.} A slave may hold SCL low; scan sampled SCL states and return
when (or whether) it releases.

\textbf{Hints.}
\begin{itemize}
\item Return the index where SCL is first high (released); $-1$ if it never
  releases (timeout).
\end{itemize}

\attempt

\textbf{Solution.}
\begin{cgym}
#include <stdio.h>

static int wait_scl_release(const int *scl_samples, int n) {
    for (int i = 0; i < n; i++)
        if (scl_samples[i]) return i;     /* SCL went high -> slave released */
    return -1;                            /* never released -> timeout */
}

int main(void) {
    int scl[] = {0, 0, 1, 1};             /* low, low, then released */
    printf("released@%d\n", wait_scl_release(scl, 4));   /* 2 */
    return 0;
}
\end{cgym}

After the master drives SCL high it must \emph{verify} SCL actually rose ---
a stretching slave holds it low until ready. The master proceeds only once it
reads SCL high, and bounds the wait with a timeout so a hung slave doesn't lock
the bus forever (on timeout, trigger the 9-clock bus-clear recovery). This
SCL-read-back is exactly why bit-banged I2C must sample SCL, not just drive it.

\followups
\begin{enumerate}
\item \emph{``What do you do on timeout?''} $\rightarrow$ Bus-clear: toggle SCL
  up to 9 times, then STOP; if still stuck, reset the slave.
\item \emph{``Does hardware I2C handle stretching?''} $\rightarrow$ Yes, it
  waits automatically, but you still set a timeout.
\item \emph{``Why might a slave stretch?''} $\rightarrow$ It needs time (ADC
  conversion, flash write) before the next byte.
\end{enumerate}
\end{gymproblem}

%---------------------------------------------------------------
\begin{gymproblem}{Choose the bus from constraints}

\textbf{Problem.} Given whether pins are scarce, throughput is high, and the
peer is asynchronous, pick UART, SPI, or I2C.

\textbf{Hints.}
\begin{itemize}
\item Async peer $\rightarrow$ UART; high speed $\rightarrow$ SPI; pin-scarce
  many-device $\rightarrow$ I2C.
\end{itemize}

\attempt

\textbf{Solution.}
\begin{cgym}
#include <stdio.h>

typedef enum { BUS_I2C, BUS_SPI, BUS_UART } bus_t;

static bus_t choose_bus(int pins_scarce, int high_speed, int async_peer) {
    if (async_peer) return BUS_UART;      /* no shared clock available */
    if (high_speed) return BUS_SPI;       /* fastest, full-duplex      */
    if (pins_scarce) return BUS_I2C;      /* two wires, many devices   */
    return BUS_SPI;                       /* default for on-board peripherals */
}

int main(void) {
    printf("%d %d %d\n",
           choose_bus(1, 0, 0),   /* I2C  -> 0 */
           choose_bus(0, 1, 0),   /* SPI  -> 1 */
           choose_bus(0, 0, 1));  /* UART -> 2 */
    return 0;
}
\end{cgym}

This encodes the selection logic an interviewer wants you to reason aloud: a
console or an off-board module with its own clock domain is UART; a fast ADC or
display is SPI; a cluster of slow sensors/EEPROMs on two wires is I2C. The
decision is about clocking model, throughput, pin budget, and device count ---
not a one-size answer.

\followups
\begin{enumerate}
\item \emph{``High speed \emph{and} pin-scarce?''} $\rightarrow$ Tension ---
  consider QSPI, or accept I2C's limit, or spend the pins on SPI.
\item \emph{``Many devices but need speed?''} $\rightarrow$ SPI with several
  chip-selects, or partition onto multiple buses.
\item \emph{``Why is UART the choice for an async peer?''} $\rightarrow$ No
  shared clock line is needed; both ends self-clock at the agreed baud.
\end{enumerate}
\end{gymproblem}

%=====================================================================
\section{Link to Her Resume}
%=====================================================================

\begin{resumelink}
\textbf{Bus selection \& bring-up:} ``On the autopilot and supporting boards I
work across SPI, I2C, UART and the differential buses --- I choose by clocking
model and pin budget, and I bring an SPI slave up by matching CPOL/CPHA from its
datasheet before anything else, scoping SCLK/MOSI/MISO at the device when reads
are garbled.''

\medskip
\textbf{I2C field recovery:} ``A stuck I2C bus usually means a slave caught
mid-byte or stretching SCL forever --- I recover with the manual 9-clock
bus-clear and a STOP, not just a peripheral re-init, and I check pull-up sizing
against bus capacitance when failures are intermittent.''

\medskip
\small Maps to the multi-protocol board bring-up and oscilloscope-based
root-cause work on the resume (autopilot boards, supporting boards, lab
instruments). Keep claims to what the resume states; deep project scripting is
Chapter 15, from \code{inputs/INTAKE.md} only.
\end{resumelink}

%=====================================================================
\section{Self-Test Scorecard}
%=====================================================================

\begin{scorecard}
\begin{enumerate}
\item Name the four SPI lines; which two carry data and in which direction?
\item What do CPOL and CPHA set, and what happens on a mode mismatch?
\item How does a master select a slave on SPI vs I2C?
\item Draw (in words) a full I2C write transaction.
\item What is clock stretching and how must the master handle it?
\item Explain I2C multi-master arbitration --- why is it non-destructive?
\item How do you size I2C pull-ups (the two bounds)?
\item Why does I2C have a speed ceiling that SPI doesn't?
\item On an I2C read, when does the master NACK and why?
\item Pick a bus: a cluster of slow sensors on two wires --- which and why?
\end{enumerate}
\end{scorecard}

\begin{answerkey}
\begin{enumerate}
\item SCLK, MOSI, MISO, CS; MOSI (master$\rightarrow$slave) and MISO
  (slave$\rightarrow$master) carry data, full-duplex.
\item CPOL = idle clock level, CPHA = sampling edge; a mismatch samples on the
  changing edge $\rightarrow$ every bit garbled/shifted.
\item SPI: assert the slave's dedicated CS. I2C: send the slave's address on the
  shared bus.
\item START $\rightarrow$ addr+W $\rightarrow$ ACK $\rightarrow$ register
  $\rightarrow$ ACK $\rightarrow$ data $\rightarrow$ ACK $\rightarrow$ STOP.
\item A slave holds SCL low to pause the master; the master must wait until SCL
  actually rises, with a timeout for a hung slave.
\item Each master reads SDA while driving; one that sends 1 but reads 0 lost and
  backs off --- the wired-AND bus leaves the winner's data intact.
\item $R_{\min}=V_{cc}/I_{\text{sink}}$ (low-side current limit);
  $R_{\max}$ from the rise-time/capacitance limit --- typically a few
  k$\Omega$.
\item Open-drain only pulls low; the rising edge is restored passively through
  the RC, so rise time bounds the clock. SPI push-pull drives both edges.
\item It NACKs the final byte to tell the slave to stop driving before STOP;
  ACKing it would make the slave send/hold more.
\item I2C --- two wires for many addressable low-speed devices, minimal pins.
\end{enumerate}
\end{answerkey}

\begin{partnote}
\emph{End of Chapter 6.} Next: \textbf{CAN Bus Deep Dive} --- frame fields,
arbitration, bit stuffing, error states, ACK, bit timing, and drone-servo
control over CAN (the lab work on the resume).
\end{partnote}
